
% ================================
% SFU CMPT 276 / Project Phase 3 Report
% ================================

\documentclass[12pt]{article}

% ---------------- Packages ----------------
\usepackage[margin=1in]{geometry}
\usepackage{setspace}
\usepackage{fancyhdr}
\usepackage{lastpage}
\usepackage{hyperref}
\usepackage{amsmath, amssymb}
\usepackage{graphicx}
\usepackage{enumitem}
\usepackage{listings}
\usepackage{xcolor}
\usepackage{float}
\usepackage{booktabs}

% ---------------- Code formatting ----------------

\definecolor{dkgreen}{rgb}{0,0.6,0}
\definecolor{mauve}{rgb}{0.58, 0.0, 0.82}
\lstset{
    language=Java,
    basicstyle=\footnotesize\ttfamily,
    numbers=left,
    frame=single,
    keywordstyle=\color{blue},
    commentstyle=\color{dkgreen},
    stringstyle=\color{mauve},
    breaklines=true,
    tabsize=4
}

% ---------------- Page Style ----------------
\pagestyle{fancy}
\fancyhf{}
\lhead{CMPT 276 -- Project Phase 3}
\rhead{Group \#15}
\cfoot{Page \thepage\ of \pageref{LastPage}}


% ---------------- Metadata ----------------
\title{\textbf{Project Phase 3: Testing}\\[0.5em]
CMPT 276: Introduction to Software Engineering\\
Simon Fraser University}

\author{
\textbf{Group \#15}\\[0.75em]
\begin{tabular}{l r}
Megan Chau      & Student \#301614135 \\
Samantha Gan    & Student \#301614044 \\
Krish Kaushik   & Student \#301628888 \\
Tanveer Muntaseer & Student \#301620006 \\
Lucas Nguyen    & Student \#301629229 \\
\end{tabular}
}

\date{March 24, 2026}

% ================================
\begin{document}
% ================================

% -------- Cover Page --------
\maketitle
\thispagestyle{empty}
\newpage

% -------- Table of Contents --------
\tableofcontents
\thispagestyle{empty}
\newpage

% Reset page numbering after TOC
\setcounter{page}{1}

% ================================
\section{Overview}
% ================================

\noindent
This report covers the testing phase (Phase 3) of \textit{Meowmino's Delivery}, a 2D arcade game built with Java and LibGDX for CMPT 276. The game features a cat on a scooter delivering orders across a 3-part town map while collecting fish and avoiding police and puddle enemies.

\medskip\noindent
We wrote \textbf{248 JUnit tests across 25 test classes}. All 248 tests pass with zero failures or errors. Tests are run automatically via Maven (\texttt{mvn test}), and code coverage is measured with JaCoCo. The test suite covers entity logic, utility classes, state machines, and multi-component integrations.

\medskip\noindent
\renewcommand{\arraystretch}{1.3}
\begin{table}[H]
\centering
\begin{tabular}{lr}
\toprule
\textbf{Metric} & \textbf{Value} \\
\midrule
Test classes & 25 \\
Total tests run & 248 \\
Failures / Errors / Skipped & 0 / 0 / 0 \\
JUnit version & 4.13.2 \\
Coverage tool & JaCoCo 0.8.11 \\
\bottomrule
\end{tabular}
\caption{Test suite summary}
\end{table}

% ================================
\section{Unit Tests}
% ================================

Unit tests verify individual features in isolation. Each test class targets a single component, and entity objects are constructed without an active OpenGL context using the \texttt{GdxTestSetup} infrastructure described in Section~\ref{sec:infrastructure}.

\subsection{Features Tested}

\renewcommand{\arraystretch}{1.3}
\begin{table}[H]
\centering
\begin{tabular}{p{0.32\textwidth} p{0.28\textwidth} p{0.30\textwidth}}
\toprule
\textbf{Feature} & \textbf{Test Class(es)} & \textbf{What Is Verified} \\
\midrule
Player movement \& collision &
  \texttt{PlayerCollisionTest} &
  Barrier pixel lookup blocks movement; position unchanged when blocked \\

\addlinespace
Bike mount / dismount &
  \texttt{MountDismountTest} &
  Q key transitions between \texttt{Player} (on-bike) and \texttt{CatOffBike}; speed and size change correctly \\

\addlinespace
Off-bike movement \& animation &
  \texttt{CatOffBikeTest} &
  Directional movement updates position; animation frame index advances per direction \\

\addlinespace
Fish collection &
  \texttt{FishTest}, \texttt{FishCollectionMathTest} &
  Fish collected within 60-unit radius; fish not collected outside radius; score increments \\

\addlinespace
Police state machine &
  \texttt{PoliceEnemyTest} &
  NONE $\to$ ALERTED $\to$ CHASING transitions; 1-second alert delay; detection range 400 units; catch threshold 75 units \\

\addlinespace
Puddle slow effect &
  \texttt{PuddleEnemyTest}, \texttt{PuddleLogicTest} &
  Player speed reduced to 70 on entry; police chase speed reduced to 120; normal speed restored on exit \\

\addlinespace
House order lifecycle &
  \texttt{HouseTest} &
  Orders activate on timer; 10-second expiry; E-key delivery awards 10 fish; delivery rejected outside 150-unit radius \\

\addlinespace
Gate passage logic &
  \texttt{GateCheckerTest}, \texttt{GateLogicTest} &
  \texttt{GateChecker} returns \texttt{CAN\_CLEAR}, \texttt{INSUFFICIENT\_FISH}, or \texttt{NOT\_AT\_GATE} based on fish count and player position \\

\addlinespace
Respawn penalty &
  \texttt{RespawnTest}, \texttt{RespawnPenaltyTest} &
  Respawn deducts 10 fish (floor 0); adds 30 seconds to timer; repositions player to map start \\

\addlinespace
Game timer &
  \texttt{TimerTest} &
  Countdown from 300 s; pause and resume; expiry fires callback at 0 \\

\addlinespace
Barrier pixel lookup &
  \texttt{BarrierLookupTest} &
  Alpha $\geq$ 128 returns blocked; alpha $< 128$ returns passable; 4-corner hit-zone tested \\

\addlinespace
Camera clamping &
  \texttt{GameCameraTest} &
  Camera clamps to map edges; follows player mid-map; \texttt{setMaxX()} and \texttt{setYBounds()} update clamp correctly \\

\addlinespace
Cinematic bars state machine &
  \texttt{CinematicBarsTest} &
  HIDDEN $\to$ APPEARING $\to$ VISIBLE $\to$ DISAPPEARING; height interpolation; mid-animation reversals \\

\addlinespace
Obama dialogue reveal &
  \texttt{ObamaDialogueHelperTest}, \texttt{ObamaSceneTest} &
  Progressive character reveal; whitehouse trigger zone; dialogue phase timing \\

\bottomrule
\end{tabular}
\caption{Unit tests: features covered}
\end{table}

% ================================
\section{Integration Tests}
% ================================

Integration tests verify interactions between two or more components. The following tests exercise cross-component flows that unit tests cannot cover in isolation.

\renewcommand{\arraystretch}{1.3}
\begin{table}[H]
\centering
\begin{tabular}{p{0.28\textwidth} p{0.28\textwidth} p{0.35\textwidth}}
\toprule
\textbf{Interaction} & \textbf{Test Class} & \textbf{What Is Verified} \\
\midrule
Police catch $\to$ Respawn &
  \texttt{PoliceRespawnIntegrationTest} &
  When \texttt{PoliceEnemy.isCatching()} is true, \texttt{Respawn.respawn()} deducts 10 fish, adds 30 s, repositions player and all police. Multiple sequential respawns each apply the penalty independently. \\

\addlinespace
Delivery $\to$ Fish award &
  \texttt{DeliveryFishIntegrationTest} &
  Completing a delivery via \texttt{House.deliver()} increments the player fish count by 10 and marks the order as inactive. \\

\addlinespace
Respawn $\to$ Timer &
  \texttt{RespawnTimerIntegrationTest} &
  Calling \texttt{Respawn.respawn()} while the timer is running correctly adds 30 s to the remaining countdown without resetting it. \\

\addlinespace
Gate check $\to$ Fish count &
  \texttt{GateFishIntegrationTest} &
  \texttt{GateChecker} reads the live fish count from the player; reducing fish below the gate threshold changes the result from \texttt{CAN\_CLEAR} to \texttt{INSUFFICIENT\_FISH} on the next call. \\

\addlinespace
Puddle entry $\to$ Respawn &
  \texttt{PuddleRespawnIntegrationTest} &
  Entering a puddle deducts 1 fish; if that deduction causes fish to drop to 0, subsequent respawn correctly applies the penalty without underflowing below 0. \\

\bottomrule
\end{tabular}
\caption{Integration tests: component interactions covered}
\end{table}

% ================================
\section{Test Infrastructure}
\label{sec:infrastructure}
% ================================

\subsection{The Challenge: Headless Testing of LibGDX Entities}

LibGDX entity constructors (e.g., \texttt{Player}, \texttt{PoliceEnemy}, \texttt{Fish}) call \texttt{new Texture(...)} which requires an active OpenGL context. JUnit tests run headlessly (no GPU or display), so these constructors cannot be called directly.

\subsection{GdxTestSetup — Shared Base Class}

All entity test classes extend \texttt{GdxTestSetup}, which:

\begin{enumerate}[itemsep=2pt]
  \item \textbf{Stubs all LibGDX singletons} (\texttt{Gdx.app}, \texttt{Gdx.files}, \texttt{Gdx.gl}, \texttt{Gdx.audio}, \texttt{Gdx.input}) via \textbf{Mockito} in a \texttt{@BeforeClass} method. This runs once per test class.
  \item \textbf{Provides factory methods} (\texttt{makePlayer}, \texttt{makeCat}, \texttt{makePolice}, \texttt{makeFish}, \texttt{makePuddle}) that use \textbf{Objenesis} to allocate entity objects without calling their constructors, then set required fields via Java reflection.
  \item \textbf{Exposes reflection helpers} (\texttt{setField}, \texttt{setEnumField}) used by subclasses to set private fields.
\end{enumerate}

\subsection{Testing Constructors (Alternative Pattern)}

Several entity classes also expose a \texttt{(float x, float y, boolean testing)} constructor that sets all \texttt{Texture} and \texttt{Sound} fields to \texttt{null}, allowing direct construction in tests that prefer it (e.g., \texttt{PoliceEnemyTest}).

\subsection{Pure-Logic Tests (No Base Class Needed)}

Tests for utility classes with no LibGDX dependency — \texttt{GateChecker}, \texttt{Respawn}, \texttt{Timer}, \texttt{ObamaDialogueHelper} — use plain JUnit with no base class or mocking. These run the fastest and have the highest coverage.

% ================================
\section{Code Coverage}
% ================================

Coverage was measured using \textbf{JaCoCo 0.8.11}, configured in \texttt{pom.xml} to run automatically during \texttt{mvn test}. The HTML report is generated at \texttt{target/site/jacoco/index.html}.

\subsection{Overall Coverage}

\renewcommand{\arraystretch}{1.3}
\begin{table}[H]
\centering
\begin{tabular}{lrrr}
\toprule
\textbf{Counter} & \textbf{Covered} & \textbf{Total} & \textbf{Coverage} \\
\midrule
Instruction & 2172 & 10428 & 20.8\% \\
Branch      & 210  & 953   & 22.0\% \\
Line        & 374  & 1826  & 20.5\% \\
Method      & 72   & 210   & 34.3\% \\
Class       & 19   & 33    & 57.6\% \\
\bottomrule
\end{tabular}
\caption{Overall JaCoCo coverage}
\end{table}

\subsection{Per-Class Line Coverage}

\renewcommand{\arraystretch}{1.3}
\begin{table}[H]
\centering
\begin{tabular}{lrrr}
\toprule
\textbf{Class} & \textbf{Lines Covered} & \textbf{Total Lines} & \textbf{Coverage} \\
\midrule
ObamaDialogueHelper & 21 & 22  & 95.5\% \\
Respawn             & 12 & 13  & 92.3\% \\
Timer               & 26 & 29  & 89.7\% \\
GameCamera          & 17 & 19  & 89.5\% \\
House               & 37 & 44  & 84.1\% \\
GateChecker         &  5 &  6  & 83.3\% \\
PoliceEnemy         & 107 & 143 & 74.8\% \\
Player              & 63 & 88  & 71.6\% \\
CinematicBars       & 21 & 39  & 53.8\% \\
Fish                &  9 & 20  & 45.0\% \\
CatOffBike          & 43 & 111 & 38.7\% \\
PuddleEnemy         &  6 & 16  & 37.5\% \\
SettingsScreen      &  1 & 34  &  2.9\% \\
ObamaScreen         &  0 & 199 &  0.0\% \\
GameScreen          &  0 & 763 &  0.0\% \\
StartScreen         &  0 & 114 &  0.0\% \\
EndScreen           &  0 & 75  &  0.0\% \\
DeliveryNotification &  0 & 38  &  0.0\% \\
InstructionsScreen  &  0 & 26  &  0.0\% \\
GameLauncher        &  0 & 12  &  0.0\% \\
Main                &  0 &  3  &  0.0\% \\
\bottomrule
\end{tabular}
\caption{Per-class line coverage (enum inner classes omitted; all show 100\%)}
\end{table}

\subsection{Why Overall Coverage Is 20.5\%}

The overall line coverage of \textbf{20.5\%} is low because the game's largest classes --- \texttt{GameScreen} (763 lines), \texttt{ObamaScreen} (199 lines), \texttt{StartScreen} (114 lines), and related screen classes --- are tightly coupled to the LibGDX \texttt{Screen} lifecycle and the OpenGL rendering pipeline. These classes cannot be exercised in a headless JUnit environment without a live display context.

Screen classes account for approximately \textbf{1215 of 1826 source lines (66.5\%)} of the codebase. Excluding them, the \textbf{logic layer} (entity classes, utility classes, state machines) has an average line coverage of \textbf{72\%} across its core classes.

\medskip
\noindent This split is by design. The classes with 0\% coverage contain rendering, input handling, and asset loading code that requires a running LibGDX application to test. The game logic that determines outcomes --- collision resolution, state machine transitions, scoring, respawn penalties, gate conditions, camera clamping --- has been extracted into testable utility and entity classes, which we have targeted for coverage.

% ================================
\section{Test Quality}
% ================================

\subsection{Measures Taken}

\begin{itemize}[itemsep=3pt]
  \item \textbf{Meaningful assertions}: Every test method contains at least one \texttt{assertEquals} or \texttt{assertTrue} that verifies observable state, not just that code runs without an exception.
  \item \textbf{Boundary conditions}: Collection-radius tests check positions exactly at the boundary (e.g., 59.9 units away and 60.1 units away from a fish) as well as well inside and well outside.
  \item \textbf{State isolation}: \texttt{@Before} methods on all entity tests reset the object under test to a known baseline before each test method, preventing cross-test contamination.
  \item \textbf{Edge cases}: The respawn tests include scenarios where the fish count is already 0 (verifying the floor), and the timer tests verify behaviour at exactly 0 seconds remaining.
  \item \textbf{Integration test realism}: \texttt{PoliceRespawnIntegrationTest} constructs both a \texttt{PoliceEnemy} and a \texttt{Player} and drives their interaction through the full catch $\to$ respawn $\to$ reposition cycle, with separate assertions on each affected field.
  \item \textbf{Regression coverage}: Tests were written alongside production code changes. Each production bug discovered during Phase 3 (see Section~\ref{sec:findings}) was accompanied by a test that fails without the fix.
\end{itemize}

\subsection{Test Organisation}

Tests are organised into three categories, following the Maven convention of \texttt{src/test/java}:

\begin{itemize}[itemsep=2pt]
  \item \textbf{Unit tests} (21 classes): one class per component, covering a single feature in isolation.
  \item \textbf{Integration tests} (5 classes with \texttt{IntegrationTest} suffix): multi-component scenarios.
  \item \textbf{Infrastructure} (1 class): \texttt{GdxTestSetup} (shared base class, not a test).
\end{itemize}

% ================================
\section{Findings}
\label{sec:findings}
% ================================

\subsection{Production Code Changes During Testing}

Writing tests required several changes to production code to make components testable and to fix discovered bugs:

\begin{enumerate}[itemsep=4pt]

  \item \textbf{Added testing constructors to entity classes.}
    \texttt{Player}, \texttt{CatOffBike}, \texttt{PoliceEnemy}, \texttt{Fish}, and \texttt{PuddleEnemy} each received a secondary constructor with a \texttt{boolean testing} flag that sets all \texttt{Texture} and \texttt{Sound} fields to \texttt{null}. This allows test code to construct entities directly without bypassing constructors via Objenesis.

  \item \textbf{Extracted \texttt{GateChecker} as a standalone utility class.}
    Gate-passage logic was previously embedded inside \texttt{GameScreen.checkGates()}. We extracted it into \texttt{GateChecker}, a pure-logic class with a \texttt{GateResult} enum (\texttt{NOT\_AT\_GATE}, \texttt{INSUFFICIENT\_FISH}, \texttt{CAN\_CLEAR}). This made gate logic independently testable and improved separation of concerns.

  \item \textbf{Extracted \texttt{ObamaDialogueHelper} from \texttt{ObamaScreen}.}
    Dialogue reveal logic and whitehouse trigger-zone checks were previously inside \texttt{ObamaScreen.render()}. Extracting them into \texttt{ObamaDialogueHelper} (a pure-Java, no-LibGDX class) enabled unit testing of all dialogue state transitions without a display.

  \item \textbf{Fixed respawn fish-count floor bug.}
    A test in \texttt{RespawnPenaltyTest} revealed that calling \texttt{Respawn.respawn()} when the player had 0 fish set the count to \texttt{-10}. The fix added a \texttt{Math.max(0, fishCount - 10)} guard. The existing test now passes and serves as a regression test.

  \item \textbf{Fixed police catch-state guard.}
    \texttt{PoliceRespawnIntegrationTest} revealed that \texttt{PoliceEnemy.isCatching()} returned \texttt{true} if the player was within 75 units even when the police was in the \texttt{NONE} or \texttt{ALERTED} state. The fix added a state guard requiring \texttt{CHASING} before proximity is checked.

\end{enumerate}

\subsection{Bugs Revealed and Fixed}

Two production bugs were uncovered and fixed by the test suite:

\begin{itemize}[itemsep=4pt]
  \item \textbf{Respawn underflow}: Fish count could go negative after respawn at 0 fish. Fixed with \texttt{Math.max(0, fishCount - 10)}.
  \item \textbf{Police catches player in wrong state}: Police could trigger a catch while still in the alert phase. Fixed by guarding \texttt{isCatching()} with a state check.
\end{itemize}

\subsection{General Quality Improvements}

\begin{itemize}[itemsep=3pt]
  \item Extracting \texttt{GateChecker} and \texttt{ObamaDialogueHelper} from screen classes improved cohesion and reduced the size of \texttt{GameScreen} and \texttt{ObamaScreen}.
  \item The testing constructors added to entity classes also served as documentation of which fields are required for an entity to function.
  \item Writing integration tests clarified the contract between \texttt{Respawn}, \texttt{Timer}, and \texttt{PoliceEnemy}, which led to tightening several implicit assumptions about call order.
\end{itemize}

% ================================
\end{document}
% ================================
